% Figure: boundary-layer growth over a flat plate. The laminar layer grows
% as sqrt(x); past transition the turbulent layer grows faster, roughly as
% x^{4/5}. The freestream U_infty is uniform above the layer. Curves are
% normalised so the laminar branch reaches a convenient height at the
% transition station.
\begin{figure}[htbp]
\centering
\begin{tikzpicture}
\begin{axis}[
  width=12.5cm, height=6.5cm,
  xlabel={streamwise distance $x$ (arb.\ units)},
  ylabel={layer thickness $\delta$},
  xmin=0, xmax=10, ymin=0, ymax=3.2,
  axis lines=left,
  legend pos=north west,
  legend cell align=left,
  grid=major, grid style={enggray!20},
  tick label style={font=\scriptsize},
  label style={font=\small},
  legend style={font=\scriptsize},
]
% Laminar: delta = 0.55 sqrt(x), drawn 0..4 (transition at x=4).
\addplot[engblue, very thick, domain=0:4, samples=120] {0.55*sqrt(x)};
\addlegendentry{laminar $\delta\sim\sqrt{x}$}
% Turbulent: delta = c x^{0.8}, matched at x=4 to laminar value 0.55*2=1.1
% so c = 1.1 / 4^0.8 = 1.1/3.031 = 0.363. Draw 4..10.
\addplot[engred, very thick, domain=4:10, samples=120] {0.363*x^(0.8)};
\addlegendentry{turbulent $\delta\sim x^{4/5}$}
% transition marker
\addplot[engamber, thick, dashed] coordinates {(4,0) (4,3.2)};
\node[font=\scriptsize, anchor=south, engamber] at (axis cs:4,3.0) {transition};
% freestream arrows
\node[font=\scriptsize, enggray] at (axis cs:1.2,2.8) {$U_\infty$};
\addplot[enggray, ->] coordinates {(0.2,2.6) (2.0,2.6)};
\end{axis}
\end{tikzpicture}
\caption[Boundary-layer growth over a flat plate]{Boundary-layer growth
over a flat plate. The laminar layer thickens as $\sqrt{x}$; downstream of
transition the turbulent layer thickens faster, roughly as $x^{4/5}$, and
is markedly thicker than a laminar layer continued to the same station.
The transition Reynolds number for a smooth plate in low-disturbance flow
is near $\mathrm{Re}_x\approx5\times10^{5}$.}
\label{fig:vol03ch12:bl-growth}
\end{figure}
